\chapter{Declarative Optimization and Equality Saturation}
\label{ch:equality-saturation}

Simplifying an expression can expose one optimization while hiding another.
Expanding $a(b+c)$ reveals separate products; factoring $ab+ac$ reveals a shared multiplier.
A destructive rewrite chooses one representation and discards the other unless it keeps an explicit history.
An equality-based optimizer can retain both and postpone selection until it knows more about the target cost.
The difficulty moves from choosing rewrite order to controlling search, proving each equality, and extracting a realizable program.

\section{A rule is a semantic claim}
\label{sec:ch62-rules}

Our runnable example uses total, pure arithmetic modulo $256$.
Every addition and multiplication returns its low eight bits, variables denote byte values, and evaluating an operand has no side effects or traps.
Under this contract, $x+0=x$, $x\cdot1=x$, and $x\cdot0=0$ hold for every input.
The last rule is not a general license to erase an effectful operand: evaluating \code{read()} and multiplying its result by zero still consumes input.

A declarative rule has a pattern, a replacement, and semantic conditions.
Those conditions may involve types, arithmetic flags, alias facts, or effect summaries.
Table~\ref{tab:ch62-legality} contrasts the byte model with richer IRs.
Rules do not acquire validity from looking algebraic.
Floating-point reassociation, for example, changes rounding and may interact with NaNs and signed zero.

\begin{table}[htbp]
\caption{Rewrite legality depends on the operation's full contract.}
\label{tab:ch62-legality}
\centering
\begin{tabular}{ll}
\toprule
Candidate & Required justification \\
\midrule
$x+0 \to x$ & correct numeric and exceptional semantics \\
$x\cdot0 \to0$ & discarded evaluation has no required behavior \\
load reuse & same bytes and compatible intervening effects \\
reassociation & exact arithmetic or an explicit relaxation \\
guard removal & assumption proved on every reaching path \\
\bottomrule
\end{tabular}
\end{table}

\section{Represent equivalence without enumerating trees}
\label{sec:ch62-structure}

An \term{e-class} represents expressions known equal under the chosen semantics.
An \term{e-node} records an operator and a tuple of child e-classes.
An \term{e-graph} stores these classes and nodes while maintaining congruence: equal arguments to the same pure operator produce equal results.
This compactly shares families of equivalent expressions, although the graph can still grow beyond any useful memory budget~\cite{willsey-2021-egg}.

Start with $(x+0)\cdot1$.
Create classes for $x$, $0$, $1$, $x+0$, and the product.
Applying the additive identity merges the class of $x+0$ with the class of $x$.
Applying the multiplicative identity merges the product with the class of its first operand.
The final root class contains a one-node representative $x$ and longer equivalent representatives.
Figure~\ref{fig:ch62-class} depicts these alternatives without claiming that each class is a distinct runtime operation.

\begin{figure}[htbp]
\centering
\begin{tikzpicture}[node distance=14mm]
\node[basic block] (root) {root e-class};
\node[basic block,below left=of root] (x) {x};
\node[basic block,below=of root] (add) {add(root, zero)};
\node[basic block,below right=of root] (mul) {mul(root, one)};
\draw[cfg edge] (root)--(x); \draw[cfg edge] (root)--(add); \draw[cfg edge] (root)--(mul);
\end{tikzpicture}
\caption{Identity rewrites put several finite expressions in one class; recursive class references also describe longer alternatives.}
\label{fig:ch62-class}
\end{figure}

Union-find tracks class identities, and a hash-consing table maps canonical operator/child tuples to their classes.
After merging child classes, two formerly distinct parent tuples may become equal.
The memo table is then stale until the implementation repairs congruence.
This repair is called \term{rebuilding}.
The companion uses repeated whole-graph scans for clarity; optimized implementations maintain affected-parent worklists and batch repairs~\cite{willsey-2021-egg}.

\begin{algorithm}[htbp]
\caption{Rebuild congruence after class unions.}
\label{alg:ch62-rebuild}
\begin{pseudocode}
\Require finite e-node list, union-find class relation
\Ensure equal canonical e-nodes have the same owning class
\Repeat
  \State clear memo table; $changed\gets false$
  \ForAll{e-nodes $n$ with owner $c$}
    \State replace each child ID in $n$ by its current representative
    \If{the canonical node already has owner $d$ in the memo table}
      \State union $c$ and $d$; record whether a new merge occurred
    \Else
      \State insert the canonical node with owner $c$
    \EndIf
  \EndFor
\Until{no new merge occurred}
\State canonicalize memo-table owners
\end{pseudocode}
\end{algorithm}

Algorithm~\ref{alg:ch62-rebuild} terminates because every productive scan reduces the number of classes.
With $C$ allocated class IDs and $N$ bounded-arity e-nodes, it needs at most $C$ scans.
Bounding each representative lookup by $O(C)$ gives a deliberately loose $O(C^2N)$ bound with expected constant-time hashing; path compression improves actual lookup behavior.
Production union-by-rank and incremental rebuilding have substantially better behavior; the demonstration is not a complexity model for \code{egg}.

\section{Saturate within a resource budget}
\label{sec:ch62-budget}

\term{Equality saturation} repeatedly matches sound rules, adds replacement nodes, and merges equivalent classes.
Stopping because no applicable rule changes the graph differs from stopping because a time, node, or iteration limit has expired.
Both can yield a correct extracted program; only the former is a fixed point of the rules actually searched.
Neither establishes that all semantic equalities have been discovered.

The companion applies constant folding and the three identities above.
It deliberately omits expansive distributivity rules so its small examples remain inspectable.
Table~\ref{tab:ch62-trace} records the identity example under unit cost per expression-tree node.
The implementation reports budget exhaustion separately from saturation and preserves a finite original representative.

\begin{table}[htbp]
\caption{The root's cheapest known tree becomes smaller as equalities are added.}
\label{tab:ch62-trace}
\centering
\begin{tabular}{lll}
\toprule
Known equality & Root representative & Tree cost \\
\midrule
initial graph & $(x+0)\cdot1$ & $5$ \\
$x+0=x$ & $x\cdot1$ & $3$ \\
$x\cdot1=x$ & $x$ & $1$ \\
\bottomrule
\end{tabular}
\end{table}

Analysis facts attached to an e-class need their own merge semantics.
A constant-value analysis must distinguish an unknown constant from a known zero.
Discovering two different byte constants in one class signals an invalid equality or a broken implementation, not a reason to choose one arbitrarily.
Conditional rules must use facts valid for every represented expression under the rule's interpretation.
Joining may-properties as though they were must-properties is a common way to introduce unsound rules.

\section{Extract a finite executable representative}
\label{sec:ch62-extraction}

Class cycles are normal: merging $x$ with $x+0$ makes a node refer to its own class.
A recursive tree traversal without a cost fixed point can loop forever.
For strictly positive additive node costs, begin with infinite cost for each class and relax a node only when all child classes have finite witnesses.
Algorithm~\ref{alg:ch62-extract} records the chosen finite tree with each improvement.
A cycle without any grounded leaf has no finite extraction.

\begin{algorithm}[htbp]
\caption{Extract a minimum-cost finite tree from a fixed e-graph.}
\label{alg:ch62-extract}
\begin{pseudocode}
\Require rebuilt graph, positive additive node costs, root class $r$
\Ensure minimum finite tree cost for $r$, or no finite witness
\State initialize every class cost to infinity
\Repeat
  \State $changed\gets false$
  \ForAll{nodes whose child classes have finite costs}
    \State $candidate\gets$ node cost plus the sum of child costs
    \If{$candidate$ is smaller than the owning class cost}
      \State store the candidate cost and its finite witness tree
      \State $changed\gets true$
    \EndIf
  \EndFor
\Until{no cost changes}
\State \Return the root's cost and witness, if finite
\end{pseudocode}
\end{algorithm}

A minimum tree need not repeat an e-class along a root-to-leaf path: the lower occurrence is itself a cheaper witness for the repeated class because costs are positive.
Thus at most $C$ levels are necessary for $C$ classes.
With bounded operator arity, at most $C+1$ relaxation scans suffice, for $O(CN)$ scalar cost work on a canonical graph.
Explicitly copying witness trees can cost much more than this, even exponentially in the number of classes when subtrees are duplicated.
An implementation should normally keep backpointers or immutable shared witnesses and charge separately for final output size.

The minimum tree is not necessarily the minimum machine program.
Shared subexpressions, register pressure, instruction combinations, and placement make machine cost non-additive.
A negative cost would invalidate the termination argument; a learned estimator can also change ranking without proving equivalence.
Figure~\ref{fig:ch62-gates} keeps the sound equality engine separate from profitability estimates.

\begin{figure}[htbp]
\centering
\begin{tikzpicture}[node distance=13mm]
\node[basic block] (rules) {verified rules and side conditions};
\node[basic block,below=of rules] (graph) {bounded equality search};
\node[basic block,below=of graph] (extract) {cost-guided finite extraction};
\node[basic block,below=of extract] (emit) {schedule, allocate, validate};
\draw[cfg edge] (rules)--(graph); \draw[cfg edge] (graph)--(extract); \draw[cfg edge] (extract)--(emit);
\end{tikzpicture}
\caption{A cost model ranks legal alternatives; it does not authorize new equalities.}
\label{fig:ch62-gates}
\end{figure}

\section{Implementation and exercises}
\label{sec:ch62-exercises}

The \code{EGraph} companion implements class union, canonicalization, rebuilding, bounded rewriting, and positive-cost extraction.
The tests compare the original and extracted identity expressions at all $256$ byte inputs, fold $256$ multiplication examples, check congruence, and reject an ungrounded cycle.
This is finite evidence for the supplied rules and implementation, not a proof of arbitrary added rules.
Chapter~\ref{ch:profile-spatial-compilation} develops the separate problem of estimating target cost.
The rewriting sequence beginning in Chapter~\ref{ch:abstract-string-rewriting} distinguishes destructive normalization, termination, confluence, and completion.
Chapter~\ref{ch:term-graph-rewriting} explains ordinary sharing and cyclic read-back; those graphs represent computations rather than equivalence classes of alternative computations.
Chapter~\ref{ch:rewrite-strategies} compares their rule-engine scheduling obligations with equality saturation.

\begin{enumerate}
\item\label{ex:ch62-wrap} Fold $250+13$ and then multiply by $13$ in byte arithmetic.
Give both folded constants and their expression-tree costs after saturation.
\item\label{ex:ch62-congruence} Classes $a$ and $b$ are merged after nodes $f(a)$ and $f(b)$ were interned separately.
Explain what rebuilding must do and why updating only the union-find parents is insufficient.
\item\label{ex:ch62-cycle} A class contains only \code{add(self,self)}.
Does it have a finite tree representative? Explain the extraction result.
\item\label{ex:ch62-effects} Show why rewriting $0\cdot f()$ to zero can change a program even when multiplication itself cannot trap.
Give an effect condition sufficient for the byte-model rule.
\end{enumerate}

\section{Further reading}
\label{sec:ch62-reading}

The \code{egg} paper explains rebuilding, e-class analyses, and the practical organization of equality saturation~\cite{willsey-2021-egg}.
The separation between program meaning and search cost also connects to the transformation-validation methods developed earlier in the book.
